% ============================================================================
% INSERT 1 -- calibration paragraph.
% GOES AFTER the paragraph ending "...building Window Sticker visualizations."
% GOES BEFORE \begin{table}[t] (Table V).
% ============================================================================

We evaluate $t_{\text{shot}}$ in Eq.~\eqref{eq:tproxy} on the shot rate measured
on \emph{ibm\_boston}, $t_{\text{nl}} = 1/2470\,\text{Hz} \approx 405\,\mu\text{s}$,
charging every shot the same regardless of depth.
Table~\ref{tab:calibration} lists it by depth alongside $t_{\text{hw}}$, the
median per-shot wall-clock duration measured on the device at that depth, and
their ratio $s_{\text{nl}} = t_{\text{nl}}/t_{\text{hw}}$.
That ratio says how much dearer a calibrated shot is than the one the device
actually delivered, and it is the factor that places the measured hardware
frontier on the same basis as the simulated curves.
Only the QPU term is rescaled. Classical training time is real wall-clock time
and does not depend on the shot rate, so it is left unscaled.


% ============================================================================
% INSERT 2 -- circuit-preparation cost model.
% GOES AFTER "...fitting an actionable prescription curve per method on the
%   training instances and evaluating it on held-out test instances. At each
%   resource level, the frontier shows the best approximation ratio achievable
%   using the different angle-setting strategies at that cost or less."
% GOES BEFORE the "Figure~\ref{fig:pareto-overlay} shows..." paragraph.
% ============================================================================

Equation~\eqref{eq:tproxy} charges only for shots, which assumes the cost of
getting a result off the device grows with the number of samples requested.
A measurement on \emph{ibm\_boston} disagrees.
We submitted single depth-$p$ QAOA circuits for a 144-node heavy-hex MaxCut
instance at $p = 1 \dots 5$ and $M \in \{10, 10^2, 10^3, 10^4\}$, averaging
three runs per setting, and the wall-clock time to return a result stayed near
$13.87$~s across the whole range.
Compilation, transfer to the control electronics and queueing dominate, and the
shots themselves are almost free at these budgets.
A second effect compounds this. COBYLA submits one circuit per objective
evaluation rather than one per iteration, and it evaluates more often than the
requested \texttt{maxiter}, so $N = 10$ costs 15 submissions and $N = 100$ costs
114, see App.~\ref{app:cost_analysis}.
Charging every submitted circuit for its own preparation gives
\begin{equation}
    T_{\text{proxy}} = t_{\text{preprocessing}}
    + n_{\text{evals}} \left( t_{\text{prep}} + M \, t_{\text{shot}} \right)
    + \left( t_{\text{prep}} + Q \, t_{\text{shot}} \right),
    \label{eq:tproxy_prep}
\end{equation}
where $t_{\text{prep}} = 13.87$~s and $n_{\text{evals}}$ is the number of
circuits COBYLA actually submits for a requested $N$. The final sampling job is
charged once for its own preparation.

A prescription is an argmax over the grid under a given cost model, so changing
the cost model changes which $(N, M, Q)$ is advised at every budget.
We therefore rebuild the frontier from the same simulated points under
Eq.~\eqref{eq:tproxy_prep} rather than shifting the published curve sideways.
No sampling is repeated, since the map from $(N, M, Q)$ to approximation ratio
is a property of the algorithm and carries no dependence on how resources are
priced.
We also cap every family at $Q \le 10^4$ under both models.
Parameter Transfer and Fixed Angles$^\dagger$ were swept to final sampling
budgets far larger than the families that optimize their angles, which is
harmless while shots dominate the cost but decides the outcome once a fixed
per-submission charge compresses the cheap end.


% ============================================================================
% INSERT 3 -- REPLACES the three paragraphs that currently begin
%   "Figure~\ref{fig:pareto-overlay} shows the resulting actionable Pareto
%    frontier for the ten instances, drawn under the two calibration..."
%   "At the smallest calibrated budgets, Fixed Angles$^\dagger$, without..."
%   "The noise-corrected calibration pushes the entire frontier to the right..."
% ============================================================================

Figure~\ref{fig:pareto-overlay} shows the resulting actionable Pareto frontier
for the ten test instances under both resource models, overlaid with the
measured real-hardware frontier.

Under Eq.~\eqref{eq:tproxy}, left panel, the training-free methods hold the
cheapest sub-second budgets, where Fixed Angles$^\dagger$ and Parameter Transfer
exchange the frontier several times, consistent with these being among the
fastest angle-setting methods in Fig.~\ref{fig:performance}.
Fixed Angles$^\star$ takes over below one second and climbs quickly, Linear Ramp
claims a narrow window either side of $1.5$~s, and Fixed Angles$^\star$ at
$p = 6$ and $p = 7$ then holds the frontier across the rest of the simulated
range.

The circuit-preparation charge, right panel, removes the cheap end of the
frontier outright.
One circuit costs $13.87$~s, so nothing is available below that, and the
sub-second regime collapses into a narrow band between $13.9$ and $17$~s in
which Parameter Transfer and Fixed Angles$^\dagger$ swap several times at almost
the same cost.
Above that band the ordering inverts relative to Eq.~\eqref{eq:tproxy}.
Parameter Transfer is training-free and submits a single circuit, so it holds
the frontier out to roughly $160$~s, Linear Ramp takes a window to about
$550$~s, and only then does Fixed Angles$^\star$ take over and climb to its
peak.

The best approximation ratio reached is essentially unchanged between the two
panels, moving by about two tenths of a percentage point, but the budget
required to get there grows by between one and two orders of magnitude.
Reaching $90\%$ takes about $3$~s under Eq.~\eqref{eq:tproxy} and about ten
minutes under Eq.~\eqref{eq:tproxy_prep}.
Angle optimization absorbs almost all of that increase, since every COBYLA
evaluation now buys a submission whether it asks for ten shots or ten thousand.
The methods that train therefore pay a fixed toll of
$n_{\text{evals}} t_{\text{prep}}$ before their first useful sample, which is
why the training-free methods hold the frontier so much further out.

The measured hardware frontier is identical in both panels.
Its resource is real elapsed wall-clock time, classical training plus QPU time
on the same per-shot basis, so it already contains whatever submission overhead
those runs incurred and charging it again would double-count.

Two caveats apply to the right panel.
The charge makes shots nearly free below $M \approx 2.5 \times 10^4$ while our
sweep stops at $M = 10^3$, so the recommended settings sit at the ceiling of the
grid and understate what a grid re-optimized for this cost model would reach.
And $t_{\text{prep}}$ describes the submission stack we measured rather than a
property of QAOA, so it will move with the software stack and with the load on
the device.
We report both models side by side for that reason instead of replacing
Eq.~\eqref{eq:tproxy}.


% ============================================================================
% INSERT 4 -- REPLACES the whole figure* environment for fig:pareto-overlay.
% ============================================================================

\begin{figure*}[t]
    \centering
    \includegraphics[width=0.98\textwidth]{figures/Res_Cost/pareto_frontier_cost_model_comparison_Q10000.pdf}
    \caption{Actionable Pareto frontier for 144-node heavy-hex MaxCut on ten
    test instances under two resource models.
    Left, shot time only, Eq.~\eqref{eq:tproxy}.
    Right, with $t_{\text{prep}} = 13.87$~s charged to every submitted circuit,
    Eq.~\eqref{eq:tproxy_prep}.
    Both panels are rebuilt from the same simulated points under their own cost
    model, with every family capped at $Q \le 10^4$ so that all of them share
    one sampling grid.
    Line colour shows which method family is the best deployable strategy at
    that resource level.
    Solid circle-marked lines are the simulated frontier and dotted
    square-marked lines the measured real-hardware frontier on
    \emph{ibm\_boston}, both on the $405\,\mu\text{s}$ per-shot basis.
    The hardware frontier is the same in both panels because its resource is
    real elapsed wall-clock time and already includes submission overhead.
    Markers sit where a strategy takes over the frontier, and whiskers show
    $\pm 1$ standard error over the ten instances.
    The shot-rate calibration is defined in the main text and tabulated by
    depth in Table~\ref{tab:calibration}.
    Methods that optimize the QAOA angles are marked with a star, e.g. Fixed
    Angles$^\star$. Methods that have some method-parameter optimization are
    indicated without any superscript. All other methods are indicated with a
    $\dagger$.}
    \label{fig:pareto-overlay}
\end{figure*}


% ============================================================================
% INSERT 5 -- REPLACES the whole Table V (tab:calibration) environment.
% The noise-corrected columns (N_CZ, sqrt(gamma), t_nc, s_nc) are gone.
% ============================================================================

\begin{table}[t]
    \centering
    \caption{Shot-rate calibration by depth $p$, measured on \texttt{ibm\_boston}.
    $t_{\text{hw}}$ is the median per-shot wall-clock duration on the device,
    $t_{\text{nl}}$ the calibrated per-shot duration used in
    Eq.~\eqref{eq:tproxy}, and $s_{\text{nl}} = t_{\text{nl}}/t_{\text{hw}}$
    the factor the shot-dependent resource is rescaled by.}
    \label{tab:calibration}
    \begin{tabular}{cccc}
        \hline
        $p$ & $t_{\text{hw}}$ ($\mu$s) & $t_{\text{nl}}$ ($\mu$s) & $s_{\text{nl}}$ \\
        \hline
        2  & 271.27 & 404.86 & 1.4925 \\
        3  & 271.27 & 404.86 & 1.4925 \\
        4  & 271.27 & 404.86 & 1.4925 \\
        5  & 280.31 & 404.86 & 1.4443 \\
        6  & 275.79 & 404.86 & 1.4680 \\
        7  & 298.39 & 404.86 & 1.3568 \\
        8  & 280.31 & 404.86 & 1.4443 \\
        9  & 271.27 & 404.86 & 1.4925 \\
        10 & 280.31 & 404.86 & 1.4443 \\
        \hline
    \end{tabular}
\end{table}
